% paper/sections/08h_pressure_summary.tex
%
% §8 章末まとめ：精度バランスと設計上のトレードオフ
% ※ 旧所在: 08c_ccd_poisson.tex（§8.2 内）→ §8 章末に移動（CHK-087: R5）
%   理由：GFM（§8.3）・BC（§8.4）を全て読んだ後に精度全体像を示す構成に変更
%% ═══════════════════════════════════════════════════════════════════════════════
\subsection{本章のまとめ：精度バランスと設計上のトレードオフ}
\label{sec:pressure_accuracy_summary}
%% ═══════════════════════════════════════════════════════════════════════════════

本章では，変密度 PPE の導出（§\ref{sec:projection_derivation}），
CCD による $\Ord{h^6}$ 離散化（§\ref{sec:ccd_poisson_matrix}），
GFM による鋭利な界面処理（§\ref{sec:gfm}），
境界条件と数値安定性（§\ref{sec:ppe_bc_stability}）を順に論じた．
以下に，本稿の数値手法の各コンポーネントの精度を整理する．

\begin{table}[htbp]
\centering
\caption{数値手法コンポーネント別の精度まとめ}
\label{tab:accuracy_summary}
\renewcommand{\arraystretch}{1.3}
\small
\begin{tabular}{llcc}
\toprule
コンポーネント & 手法 & 空間精度 & 時間精度 \\
\midrule
CLS 移流（界面追跡） & Dissipative CCD + TVD-RK3 & $O(h^5)$ & $O(\Delta t^3)$ \\
幾何量計算（曲率等） & CCD 法 & $O(h^6)$ & — \\
PPE（離散化＋求解） & CCD 微分形式 + GFM ジャンプ補正 & $O(h^6)$ & — \\
（参考：FVM + BiCGSTAB） & （FVM + 調和平均係数） & （$O(h^2)$） & — \\
NS 粘性項 & CCD + Crank--Nicolson & $O(h^6)$ & $O(\Delta t^2)$ \\
NS 対流項（Predictor） & CCD（1次微分）+ AB2 外挿 & $O(h^6)$ & $O(\Delta t^2)$ \\
Projection スプリッティング & IPC 法（van Kan 1986） & — & $O(\Delta t^2)$ \\
GFM 界面処理（表面張力） & Young--Laplace ジャンプ直接組み込み & $O(h^2)$（界面近傍†） & — \\
\midrule
\textbf{全体（空間律速：GFM 界面処理†）} & & $\boldsymbol{O(h^2)}$（界面近傍） & — \\
\textbf{全体（時間律速：AB2+IPC 整合）} & & — & $\boldsymbol{O(\Delta t^2)}$ \\
\bottomrule
\end{tabular}

\smallskip
\noindent†\ GFM は CSF の $O(\varepsilon^2)$ 物理誤差床を除去するが，界面直近（1–2格子幅以内）では
Extension PDE なしに $O(h^6)$ 精度を完全に維持することは困難である（§\ref{sec:gfm}，今後の課題）．
平滑領域では全コンポーネント $O(h^5)$ 以上が確保される．
\end{table}

\begin{tcolorbox}[warnbox, title={精度ミスマッチと設計上のトレードオフ}]
\begin{enumerate}
  \item \textbf{空間精度：}
        全数値コンポーネントが $O(h^5)$–$O(h^6)$ に揃っている（平滑領域）．
        GFM により旧 CSF モデルの誤差床 $O(\varepsilon^2) \approx O(h^2)$ は除去されたが，
        界面直近では GFM 処理精度（$O(h^2)$ 以上）が律速段階となる．
        FVM + BiCGSTAB では $O(h^2)$ に留まる（表~\ref{tab:ccd_vs_fvm}）．

  \item \textbf{時間精度：}
        AB2 + IPC により全体 $O(\Delta t^2)$（§\ref{sec:ipc_derivation}）．
        対流（AB2），粘性（Crank--Nicolson），Projection（IPC）の3要素が統一されている．
\end{enumerate}
\end{tcolorbox}
